% ================================================================
% SECTION 5 — DISCUSSION
% Target: ~5 pages
% ================================================================

\section{Discussion}
\label{sec:discussion}

This section interprets the empirical findings in light of the four
research questions stated in Section~\ref{sec:rq}, situates the
results within the broader literature, and draws out the practical
implications for infrastructure resilience policy and freight rail
planning in South Africa.

% ---------------------------------------------------------------
\subsection{Answering the Research Questions}
\label{sec:disc_rq}
% ---------------------------------------------------------------

\textbf{RQ1: What is the causal effect of the April 2022 KZN floods
on monthly NATCOR rail freight volume, and how does this estimate
vary across five SCM-family estimators?}

The five SCM-family estimators, applied to the full five-corridor
donor pool, produce a tightly clustered set of estimates.
The average monthly treatment effect (ATT) ranges from
$-0.154$\,MT/month (MC-SCM) to $-0.159$\,MT/month (SDID), with a
consensus mean of $-0.157$\,MT/month.
Over the 33-month post-flood window, this translates to a cumulative
volume loss of approximately $-5.08$ to $-5.24$\,MT, with a
consensus mean of $-5.17$\,MT.

The across-estimator consistency is itself a substantive finding.
The five methods differ in their identification assumptions, their
treatment of donor weights, and their regularisation strategies.
Base SCM imposes a convexity constraint; ASCM relaxes it via
ridge bias-correction; EN-SCM further relaxes it through sparse
penalisation; MC-SCM abandons explicit weighting in favour of
low-rank matrix completion; SDID reweights both units and time
periods simultaneously.
The fact that all five converge within a 0.15\,MT band suggests
that the estimated effect is not an artefact of any particular
modelling assumption, but reflects a genuine and robust signal in
the data.

\textbf{RQ2: How sensitive are SCM estimates to the size and
composition of the donor pool, and what is the minimum donor pool
required for consistent counterfactual estimation?}

The comparison between the 1-donor and 5-donor configurations
provides a sharp empirical answer.
With a single donor (CAPE), the five methods produce estimates
spanning nearly 15\,MT, from $-11.1$\,MT to $+3.7$\,MT.
With five donors, the span collapses to 0.15\,MT.
This is not a gradual improvement: it is a qualitative transition
from incoherence to consensus.

The minimum donor pool required for consistent estimation cannot be
determined from this study alone, as only two configurations were
tested.
However, the results are consistent with the theoretical requirement
that the donor pool must be sufficiently rich to span the
pre-treatment trajectory of the treated unit.
CAPE alone cannot reproduce NATCOR's pre-period dynamics
(RMSPE $\approx 0.49$\,MT), while five donors reduce RMSPE to
0.026\,MT --- a 95\% improvement in fit that directly translates
into estimator coherence.
The practical implication is that researchers applying SCM to
corridor-level freight data should treat a pre-period RMSPE below
5\% of the treated unit's mean as a minimum threshold for
proceeding to causal interpretation.

\textbf{RQ3: How robust are the causal estimates to perturbations
in the pre-period window, donor pool composition, and model
hyperparameters?}

Robustness is demonstrated across three dimensions.
First, in-time placebo tests show that zero of five pseudo-treatment
dates produces an effect as large as the true estimate under any
of the five methods --- a result consistent across all structural
variants of the estimator.
Second, leave-one-out donor tests show that removing any single
corridor from the donor pool shifts the consensus estimate by at
most 0.15\,MT (SDID, LOO spread), and by as little as 0.036\,MT
(ASCM), confirming that no single corridor is driving the result.
Third, the tight agreement between the jackknife confidence interval
for SDID ($[-5.393, -5.080]$\,MT) and the block bootstrap interval
for MC-SCM ($[-5.210, -4.948]$\,MT) confirms that the precision of
the estimate is not sensitive to the choice of inference procedure.

Taken together, these three robustness checks support the conclusion
that the estimated effect is stable and not an artefact of any
particular modelling choice.

\textbf{RQ4: What is the cumulative freight volume loss attributable
to the flood shock, and what does this imply for the economic cost
of infrastructure disruptions on the NATCOR corridor?}

The consensus cumulative loss of approximately $-5.17$\,MT over 33
months represents a 4.0\% sustained reduction in NATCOR throughput
relative to its pre-flood mean.
At a conservative rail freight revenue rate of ZAR\,250 per tonne,
this implies a direct revenue loss to Transnet of approximately
ZAR\,1.28 billion.
This figure covers only the loss in freight revenue on the NATCOR
corridor itself; it does not capture second-order costs including
road infrastructure damage from diverted heavy trucks, shipper
inventory holding costs, port congestion externalities at Durban,
or the economic loss to exporters whose cargo was delayed or rerouted.
A comprehensive social cost estimate would be substantially higher.

% ---------------------------------------------------------------
\subsection{Interpretation of the Persistence Effect}
\label{sec:disc_persist}
% ---------------------------------------------------------------

Perhaps the most policy-relevant finding of this study is not the
magnitude of the immediate disruption but the persistence of the
gap.
The monthly gap series remains uniformly negative across all 33
post-flood months, with no evidence of mean reversion to the
synthetic counterfactual by December 2024.
This pattern is inconsistent with a purely physical disruption
hypothesis --- the view that rail volumes fall during the repair
period and recover fully once infrastructure is restored.
Physical repairs to the NATCOR mainline were largely completed
within six to nine months of the flood, yet the gap persists well
beyond this window.

Two complementary mechanisms are consistent with this persistence.
The first is demand-side hysteresis: shippers who diverted freight
to road transport during the disruption period may have renegotiated
contracts, invested in truck fleets, or developed road-based
logistics routines that are costly to reverse.
In freight markets characterised by switching costs and long-term
contracts, a temporary disruption can trigger permanent modal
diversion if it coincides with a contract renewal cycle.
The second mechanism is supply-side degradation: the flood may
have exacerbated pre-existing deterioration in locomotive
availability and track quality on NATCOR, reducing service
reliability in ways that deterred returning shippers even after
physical reconstruction was nominally complete.

These mechanisms are not mutually exclusive and are likely to have
operated simultaneously.
Disentangling their relative contributions would require
shipper-level survey data or commodity-specific volume breakdowns
that are not available in the aggregate corridor panel used in
this study, and represents an important direction for future
research.

% ---------------------------------------------------------------
\subsection{Relationship to Prior Literature}
\label{sec:disc_lit}
% ---------------------------------------------------------------

The finding of a causal, persistent negative effect of the 2022
KZN floods on NATCOR volumes is consistent with the broader
literature on infrastructure disruptions and freight demand.
Studies of natural disaster effects on transport networks
consistently find that recovery is slower than physical repair
timelines suggest~\parencite{cavallo2013}, and that the demand-side
response to disruption is asymmetric: modal diversion away from
rail is more rapid than modal return after service restoration.

The donor pool inadequacy result has a direct parallel in the
comparative case studies literature.
\textcite{abadie2003} warn against applying SCM when the number
of available control units is too small to approximate the treated
unit's pre-period trajectory~\parencite{abadie2003}.
The 1-donor versus 5-donor comparison in this thesis provides a
rare empirical demonstration of this failure mode in a logistics
context: the theoretical warning becomes visible as a 15\,MT
spread in cumulative effect estimates, rendering any single estimate
from the 1-donor configuration literally meaningless for policy
purposes.

The methodological contribution of systematically comparing five
SCM estimators on the same dataset fills a gap identified by
\textcite{benmichael2021}, who argue that applied researchers
should report multiple estimators as a robustness check rather
than relying on a single method~\parencite{benmichael2021}.
This thesis operationalises that recommendation in a transport
economics setting and demonstrates that, given an adequate donor
pool, the choice of estimator matters little for the point estimate
--- though it matters considerably for the inference procedure and
the interpretation of uncertainty.

% ---------------------------------------------------------------
\subsection{Limitations}
\label{sec:disc_limit}
% ---------------------------------------------------------------

Several limitations of this study should be acknowledged.

First, the donor pool is constrained to five South African rail
corridors operated by Transnet Freight Rail.
An international donor pool --- incorporating freight corridor data
from comparable emerging-market rail operators --- could in
principle improve pre-period fit further and provide stronger
identification.
However, comparable monthly corridor-level data are not publicly
available for other African rail operators, making this extension
infeasible with current data sources.

Second, the analysis uses aggregate monthly corridor volume
(million tonnes) as the outcome variable.
This measure conflates commodity types, distance bands, and
contract structures that may respond differently to the flood shock.
A more granular analysis using commodity-level or origin--destination
pair data could reveal heterogeneity in the flood effect that is
obscured in the aggregate.

Third, the SCM framework estimates the average treatment effect
over the full post-period.
It does not directly model the time profile of recovery or estimate
whether the effect is accelerating, decelerating, or constant.
Dynamic treatment effect models --- such as local projections or
state-space approaches --- could complement the SCM results by
characterising the recovery trajectory more precisely.

Fourth, the causal identification relies on the parallel trends
assumption: that in the absence of the flood, NATCOR would have
evolved as a weighted combination of the donor corridors.
This assumption is supported by the tight pre-period fit but cannot
be directly tested post-treatment.
Concurrent shocks that affected NATCOR differentially from the
donor pool --- such as a corridor-specific locomotive crisis or
a major customer insolvency --- could confound the estimated flood
effect.
The available evidence from publicly reported operational data
does not suggest such a confound, but it cannot be ruled out
with certainty.

% ---------------------------------------------------------------
\subsection{Policy Implications}
\label{sec:disc_policy}
% ---------------------------------------------------------------

The results of this study carry several concrete implications for
infrastructure resilience planning in South Africa.

The persistence of the volume gap beyond the physical repair period
suggests that resilience investment should not be evaluated solely
on the speed of infrastructure reconstruction.
A corridor that is physically repaired but commercially disrupted
--- because shippers have permanently diverted to road --- may
never recover its pre-shock throughput.
Resilience policy therefore needs to address demand-side recovery
as well as supply-side restoration, for instance through temporary
tariff incentives to retain or attract returning shippers during
the post-repair period.

The estimated revenue loss of ZAR\,1.28 billion over 33 months
provides a lower bound on the economic case for pre-emptive
resilience investment on the NATCOR corridor.
Infrastructure hardening measures --- such as flood-proofing of
critical embankments, redundant signalling systems, and rapid
deployment capabilities for bridge repair --- that cost less than
this figure over their operational lifetime would pass a simple
cost-benefit test against the observed flood disruption alone,
without counting future flood events or the avoided secondary
costs of road diversion.

Finally, the donor pool inadequacy finding has a methodological
implication for the use of causal inference in transport policy
evaluation.
Regulatory bodies and infrastructure planners who commission
SCM-based impact evaluations of corridor disruptions must ensure
that the donor pool is adequate before accepting any single-method
result.
A pre-period RMSPE that is large relative to the treated unit's
mean, or a situation in which different SCM methods produce
qualitatively different results, should be treated as a signal
that the donor pool is inadequate --- not as a reason to prefer
one method over another.
